\chapter{Pesnik}

\editorialnote{Martin Silenus, njegov život, umetnost, propadanje i Spev o Hiperionu.}

Martin Silenus je nekoliko trenutaka ćutao pre nego što je počeo.

Za čoveka koji je gotovo neprestano govorio, psovao, pio ili recitovao tuđe stihove, ta tišina bila je neobična.

Onda je podigao čašu.

``Na početku je bila Reč'', rekao je. ``Onda smo izmislili mašine koje su pisale umesto nas. Posle toga mašine koje su mislile umesto nas. I tu je, otprilike, počela smrt književnosti.''

Bron Lamija je zakolutala očima.

Silenus se osmehnuo.

``Pitali ste zašto idem Šrajku. Da biste to razumeli, morate prvo da shvatite jednu stvar. Ja nisam oduvek bio ovo.''

Pokazao je na sebe.

``Bio sam mnogo gori.''

\bigskip

Martin Silenus rođen je na Staroj Zemlji.

Pravoj Zemlji.

Ne na nekoj koloniji koja je kasnije uzela njeno ime, nego na matičnom svetu čovečanstva, pre njegovog konačnog uništenja.

Rođen je u jednoj od poslednjih veoma bogatih porodica.

Detinjstvo mu je proteklo na ogromnom porodičnom imanju, među šumama, vrtovima, slugama, androidima i ljudima koji nikada nisu morali da razmišljaju o tome koliko nešto košta.

Zemlja je već umirala.

Decenijama ranije jedna ljudska greška stvorila je malu crnu rupu u njenoj unutrašnjosti. Planeta je od tada prolazila kroz sve češće geološke katastrofe.

Silenus je odrastao gledajući jedan od najlepših svetova koje je čovek poznavao kako polako nestaje.

Možda je upravo zbog toga, govorio je, od detinjstva znao da će biti pesnik.

Nije smatrao da je izabrao poeziju.

Smatrao je da je ona izabrala njega.

Njegov učitelj, don Baltazar, terao ga je da pamti stare pesnike. Homer, Šekspir, Kits, Bajron, Jejts i desetine drugih postali su mu gotovo druga porodica.

Dok su druga deca sve više učila pomoću implantata i direktnog pristupa podacima, Silenus je učio reči.

Hiljade stihova.

Ritam.

Metar.

Zvuk jezika.

I već tada je bio potpuno siguran da je genije.

Problem je bio samo u tome što njegova poezija nije bila naročito dobra.

``Naravno'', rekao je Silenus hodočasnicima, ``ja to tada nisam znao. Loš pisac je obično poslednja osoba koja sazna da je loš pisac.''

Objavljivao je zahvaljujući porodičnim vezama.

Niko mu nije govorio istinu.

A onda je Zemlja umrla.

Sa njom je umrla i njegova majka.

Silenus je preživeo samo zato što ga je neposredno pre katastrofe poslala na putovanje sporije od svetlosti.

Put je trajao vekovima objektivnog vremena.

Za njega, u kriogenom snu, mnogo kraće.

Kada se probudio, svet koji je poznavao više nije postojao.

Zemlje nije bilo.

Porodice nije bilo.

Bogatstva gotovo da nije bilo.

A najgore od svega, dug kriogeni san oštetio mu je mozak.

Doživeo je težak moždani udar.

Mogao je da razmišlja.

Mogao je da razume ljude.

Ali gotovo nije mogao da govori.

Čitav njegov aktivni rečnik sveo se na nekoliko psovki i najjednostavnijih reči.

Za čoveka koji je verovao da je rođen zbog jezika, to je bilo gore od smrti.

\bigskip

Probudio se na Rajskoj Kapiji.

Ime je bilo gotovo uvredljivo.

Planeta je bila siromašna, prljava i jedva pogodna za život. Ljudi su bili sabijeni u malom području oko postrojenja koje je održavalo atmosferu dovoljno dobrom za disanje.

Silenus, nekada naslednik ogromnog imanja na Staroj Zemlji, završio je kao fizički radnik.

Kopao je kanale.

Čistio otpad.

Radio u blatu i otrovnom vazduhu.

Živeo među ljudima koji bi za njegov nekadašnji život bili gotovo nevidljivi.

I upravo tamo je, prema sopstvenim rečima, prvi put zaista postao pesnik.

Ranije je pisao o životu.

Na Rajskoj Kapiji počeo je da ga živi.

Radio je po čitav dan.

Noću je pokušavao da piše.

U početku bez reči.

Ideje su postojale u glavi, ali nije mogao da ih izrazi.

Onda su počele da se vraćaju.

Jedna reč.

Pa druga.

Čitave rečenice.

Mozak je polako pronalazio nove puteve oko oštećenja.

Silenus je počeo da shvata koliko je svaka reč dragocena upravo zato što je morao da je osvaja.

Ranije bi napisao prvu približno dobru reč koja mu padne na pamet.

Sada je mogao satima da čeka pravu.

``Tada sam konačno shvatio'', rekao je, ``da između gotovo prave reči i prave reči postoji ista razlika kao između gledanja munje i toga da te munja pogodi.''

Na jeftinom papiru, u straćari na Rajskoj Kapiji, nastalo je njegovo prvo veliko delo.

Nije još bilo dovršeno.

Ali bilo je dovoljno dobro da mu promeni život.

\bigskip

Do izdavača je stiglo gotovo slučajno.

Jednog dana lokalni nasilnik i njegovi prijatelji pokušali su da mu uzmu ono malo novca što je imao.

Silenus nije imao novca.

Ali imao je rukopis.

Kada je odbio da ga preda, pretukli su ga i bacili stranice u blato.

Slučajno je naišla žena lokalnog upravnika.

Videla je šta se dogodilo.

Odvezla je Silenusa u bolnicu.

Pokupila je rukopis.

I pročitala ga.

Dopao joj se dovoljno da ga preko poznanika pošalje velikoj izdavačkoj kući u Mreži.

Dok se Silenus još oporavljao od prebijanja, stigao je ugovor.

Nekoliko nedelja kasnije njegova knjiga bila je objavljena.

Urednici su je nazvali \emph{Umiruća Zemlja}.

Nisu objavili ceo rukopis.

Izbacili su filozofiju.

Eksperimente sa jezikom.

Teže delove.

Ostavili su nostalgiju za Starom Zemljom, majku, imanje, poslednje dane planete i dovoljno tuge da se čitaoci osećaju kao da su pročitali nešto važno.

Silenus je mrzeo ono što su napravili od njegove knjige.

Publika ju je obožavala.

Prodate su milijarde primeraka.

Čovek koji je nekoliko meseci ranije kopao otpad ponovo je postao bogat.

Ovoga puta bio je i slavan.

\bigskip

Za kratko vreme Martin Silenus postao je jedna od najpoznatijih ličnosti Hegemonije.

Imao je kuću čije su se sobe nalazile na različitim planetama, povezane dalekobacačima.

Iz jedne sobe mogao je da izađe na svet sa okeanom, iz druge na planinsku liticu, iz treće u centar nekog grada.

Imao je novac.

Droge.

Prijeme.

Ljubavnike.

Poznanstva sa političarima, umetnicima i slavnim ličnostima.

Imao je sve što je nekada smatrao dokazom uspeha.

I postepeno je prestao da piše.

Pokušao je da popuni prazninu drugim stvarima.

Politikom.

Religijom.

Putovanjima.

Drogama.

Seksom.

Ničim nije uspeo.

Njegov brak se raspao.

Novac je nestajao.

A onda je ponovo počeo da piše.

Ovoga puta nije želeo bestseler.

Želeo je da napiše nešto dobro.

Radio je gotovo dve godine.

Kada je završio, njegova urednica je pročitala rukopis i zaplakala.

``Remek-delo'', rekla mu je.

Silenus ju je pitao:

``Hoće li se prodavati?''

``Ne.''

Bila je u pravu.

Knjiga je doživela katastrofu.

Gotovo niko je nije kupio.

Kritičari su je nazivali pretencioznom, nerazumljivom i zastarelom.

Jedan poznati komentator priznao je da nije mogao da je čita.

Silenus je izgubio ogroman deo preostalog novca.

Izdavač mu je tada ponudio izbor.

Može da nastavi da bude pesnik i bankrotira.

Ili može ponovo da piše ono što ljudi kupuju.

Izabrao je novac.

\bigskip

Napisao je \emph{Umiruću Zemlju II}.

Ovoga puta u prozi.

Knjiga je bila napravljena prema reakcijama test-publike.

Dužina poglavlja bila je prilagođena pažnji prosečnog čitaoca.

Likovi su napravljeni tako da budu dopadljivi.

Romansa je ubačena kada su testovi pokazali da nedostaje romanse.

Avantura kada je nedostajalo uzbuđenja.

Seks kada je pažnja počela da pada.

Prodala se odlično.

Usledila je \emph{Umiruća Zemlja III}.

Onda četvrti deo.

Peti.

Šesti.

Na kraju čitav niz.

Silenus je postao ono što je najviše prezirao.

Profesionalni proizvođač sadržaja.

``Svaki put kad bih napisao još jednu'', rekao je, ``neko bi mi platio dovoljno da napišem sledeću. To je najefikasniji oblik pakla koji je izdavačka industrija ikada izmislila.''

Pio je sve više.

Pisao sve gore.

Njegovo telo ostajalo je mlado zahvaljujući tretmanima za produženje života, ali osećao je da se njegov pravi život troši.

A onda je upoznao Tužnog Kralja Bilija.

\bigskip

Vilijam XXIII, poznat svima kao Tužni Kralj Bili, bio je ekscentrični vladar malog kraljevstva u izbeglištvu.

Izgledao je nespretno.

Govorio je sa jakim mucanjem.

Oblačio se katastrofalno.

Ali posedovao je nešto što je Silenus retko sretao.

Razumeo je umetnost.

Kada su se prvi put upoznali, Bili mu je rekao da je pročitao \emph{Umiruću Zemlju}.

``Zanimljiva je'', rekao je.

Silenus mu se zahvalio.

``Zanimljiva je zato što je neko očigledno izbacio sve što je u originalu valjalo.''

Silenus je tada odlučio da mu se kralj dopada.

Bili je pročitao i Silenusovu neuspešnu poemu.

Za razliku od gotovo svih drugih, shvatio ju je.

Smatrao ju je najboljom poezijom objavljenom u Mreži generacijama.

Pozvao je Silenusa da živi među umetnicima koje je okupljao oko sebe.

Silenus je pristao.

Godinama je bio njegov štićenik.

Zatim savetnik.

Na kraju prijatelj.

Ali Silenus i dalje nije mogao da napiše ono što je želeo.

Nedostajala mu je muza.

Jednom su on i Bili dugo razgovarali o najvećim pesnicima prošlosti.

Bili ga je pitao ko je bio najčistiji pesnik koji je ikada živeo.

Silenus je razmišljao.

``Džon Kits.''

Objasnio mu je kako je Kits mlad umro, ostavivši nedovršene velike poeme i život gotovo potpuno posvećen umetnosti.

Bili je zapamtio odgovor.

Godinama kasnije pozvao je Silenusa u svoju odaju.

Pokazao mu je hologram jednog udaljenog sveta.

``Hiperion.''

Silenus je odmah prepoznao ime.

Bili mu je pokazao Vremenske grobnice.

Gradove koji su tek počinjali da se razvijaju.

Ogromne neistražene oblasti.

Svet bez duboke povezanosti sa Mrežom i bez neprekidnog prisustva Svestvari.

``Misterija'', rekao je Kralj Bili. ``Umetnicima je još potrebna misterija.''

Planirao je da tamo preseli svoju umetničku koloniju.

Silenus nije bio naročito oduševljen.

Onda ga je Bili pitao:

``Pišeš li?''

``Ne.''

``Vratila se muza?''

``Ne.''

Bili se nasmešio.

``Onda pakuj stvari. Idemo na Hiperion.''

\bigskip

Na Hiperion je stiglo više hiljada umetnika, pisaca, muzičara, arhitekata i svih ljudi koji su želeli da veruju da još mogu da naprave nešto novo.

Izgradili su gradove.

Najambiciozniji među njima bio je Grad Pesnika.

Beli tornjevi.

Galerije.

Pozorišta.

Vrtovi.

Radionice.

Mesto zamišljeno kao utočište od kulturne učmalosti Hegemonije.

Silenus je ostao tamo.

I nije napisao ništa.

Godine su prolazile.

Ostali umetnici stvarali su, svađali se, zaljubljivali i proglašavali jedni druge genijima.

Silenus je sedeo pred praznim stranicama.

Njegova tehnika bila je savršena.

Sve što bi napisao zvučalo je kao poezija.

Ali nije imalo života.

Posle gotovo deset godina odlučio je da se ubije.

Pre toga je, kako je rekao, poželeo da se malo zabavi.

Kod biovajara je promenio telo.

Dobio je dlakave noge, papke i izgled mitološkog satira.

Ostatak modifikacija napravljen je tako da odgovara reputaciji satira.

Pio je.

Spavao sa svima koji su to želeli.

Pravio skandale.

Postao jedna od lokalnih znamenitosti.

``Bilo je divno'', rekao je Silenus.

Zastao je.

``Bilo je užasno.''

A onda, upravo kada je odlučio da će se ubiti, ljudi su počeli da nestaju.

\bigskip

U početku niko nije bio naročito zabrinut.

Grad Pesnika bio je pun ekscentrika.

Ljudi bi otišli na nekoliko dana bez objašnjenja.

Onda je pronađeno prvo telo.

Pa drugo.

Neki ljudi jednostavno su nestali.

Drugi su pronađeni rastrgnuti.

Nije postojao obrazac.

Nije postojalo poznato oružje koje bi moglo da napravi takve povrede.

Grad je postavio senzore.

Formirana je policija.

Pretražene su Vremenske grobnice.

Ništa.

Ubistva su se nastavila.

I upravo tada Martin Silenus ponovo je počeo da piše.

Posle godina tišine reči su se vratile.

Ne polako.

Ne uz napor.

Tekle su.

Prvi put posle Rajske Kapije Silenus je osećao ono zbog čega je verovao da postoji.

Počeo je novu poemu.

U početku nije želeo da razmišlja o vezi između ubistava i svog iznenadnog nadahnuća.

Kralj Bili jeste.

Jednog dana pozvao ga je da pogleda snimak jednog od napada.

Kamera je slučajno zabeležila smrt žene u njenoj sobi.

Silenus je gledao kako nešto gotovo nevidljivo raskida ljudsko telo.

Na trenutak se stvorenje dovoljno dugo našlo pred kamerom.

Snimak je zaustavljen.

Uvećan.

Silenus je prvi put video lice Šrajka.

Više od tri metra visoko stvorenje.

Metal.

Hrom.

Sečiva.

Trnje.

Crvene oči.

Nešto između mašine, životinje i košmara.

``Šrajk'', rekao je.

Kralj Bili ga je pogledao.

Znao je da je Silenus istraživao stare legende o njemu još mesecima pre prvih ubistava.

``Zašto?''

``Zbog poeme.''

Silenus je objasnio šta je pronašao.

Prema kultu koji je već postojao na Hiperionu, Šrajk je Bog bola.

Biće koje dolazi izvan vremena.

Vesnik konačnog uništenja.

Neki su verovali da ga je, neposredno ili posredno, stvorilo samo čovečanstvo.

Bili ga je pitao:

``Šta je on zapravo?''

Silenus je dugo ćutao.

Onda mu je sinulo.

``Moja muza.''

\bigskip

Od tog trenutka više nije mogao da odbaci vezu.

Ubistva su počela.

On je počeo da piše.

Kada bi Šrajk ponovo ubio, Silenus bi sledećih dana radio kao opsednut.

Naslov dela promenio je u:

\emph{Spev o Hiperionu}.

Nije to bila poema o jednoj planeti.

Bila je o čovečanstvu.

O vrsti koja je uništila sopstveni matični svet, proširila se među zvezde i svuda sa sobom ponela istu oholost.

O civilizaciji koja je verovala da razume sve zato što je napravila dovoljno moćne mašine.

O kraju ljudskog doba.

Šrajk je u poemi postao odgovor na tu oholost.

Možda kazna.

Možda posledica.

Možda nešto što će čovek tek napraviti u budućnosti.

Silenus nije znao.

Ali je pisao.

I bio je siguran da nikada u životu nije pisao bolje.

To ga je užasavalo.

I ushićivalo.

Grad se istovremeno praznio.

Ljudi su odlazili.

Kralj Bili konačno je naredio evakuaciju.

Njegov san o umetničkoj utopiji bio je mrtav.

Silenus nije otišao.

Nekoliko stotina drugih takođe je ostalo.

Ubistva su se nastavila.

Broj stanovnika se smanjivao.

A Silenus je nastavljao da piše.

Godinama.

\bigskip

Grad Pesnika se raspadao oko njega.

Akvadukti su prestali da rade.

Vrtovi su se osušili.

Pustinja je polako gutala ulice.

Galerije su zarasle.

Kupole su pucale.

Ljudi su nestajali jedan po jedan.

Silenus je ostajao.

Pisao je.

Ponekad bi noću čuo zvuk metala na kamenu.

Ponekad bi bio siguran da ga nešto posmatra.

Šrajka gotovo nikada nije video.

Ali osećao je njegovo prisustvo.

Sa svakim novim ubistvom poema je rasla.

Silenus je počeo da veruje u nešto što je i sam smatrao ludilom.

Možda nije samo Šrajk stvarao poemu.

Možda je i poema stvarala Šrajka.

U početku je bila Reč.

Šta ako umetnost nije samo opis stvarnosti?

Šta ako je, dovoljno snažna, može prizvati?

Šta ako je pesnik video ono što mora da postoji, a ne samo ono što već postoji?

Silenus više nije znao da li piše proročanstvo ili ga izaziva.

Na kraju je u Gradu Pesnika ostao gotovo sam.

Tada je prestao da piše.

Nije želeo.

Jednostavno više nije mogao.

Šrajk je prestao da ubija.

Njegova muza je zaćutala.

\bigskip

Jedne noći Silenus se vratio iz šetnje do Vremenskih grobnica i pronašao nekoga u svojoj radnoj sobi.

Kralja Bilija.

Bili je bio star.

Mnogo stariji nego kada su zajedno došli na Hiperion.

Pred njim su ležale hiljade stranica \emph{Speva o Hiperionu}.

Čitao ih je.

``Impresivno'', rekao je.

Silenus je bio besan.

Niko nije imao pravo da čita nedovršeni Spev.

Ni njegov prijatelj.

Ni čovek koji mu je omogućio da dođe na Hiperion.

Bili je pokazao poslednje stranice.

``Mesecima nisi ništa napisao.''

``Pa?''

``Poslednje što si napisao nastalo je kada je umro poslednji stanovnik Grada Pesnika.''

Silenus nije odgovorio.

Bili je nastavio.

``Ne možeš da pišeš bez njega.''

``Sereš.''

``Ne možeš da dovršiš poemu dok tvoja muza ponovo ne prolije krv.''

Silenus je želeo da ga izbaci iz sobe.

Bili je izvukao nervni ošamućivač.

Silenus je pokušao da mu priđe.

Sledeće čega se sećao bilo je kameno dvorište.

Nije mogao da pomera telo.

Kralj Bili stajao je pored stare fontane.

U njoj je bio rukopis.

Bilijeve oči bile su pune suza.

U ruci je držao upaljač.

``Žao mi je, Martine.''

Zapalio je prve stranice.

Silenus je pokušao da ustane.

Nije mogao.

``Ne!''

Još stranica nestalo je u vatri.

Godine njegovog rada.

Najbolje rečenice koje je ikada napisao.

Jedino delo za koje je verovao da opravdava njegov život.

``Molim te.''

Bili nije stao.

``Ovo mora da se završi.''

Silenus je već uspevao da pomera ruke i noge.

Pokušavao je da ustane.

Tada su shvatili da više nisu sami.

\bigskip

Šrajk je stajao u dvorištu.

Niko ga nije video kako dolazi.

Kao da je jednostavno odlučio da odjednom bude tamo.

Visok.

Metalni.

Četvororuk.

Prekriven sečivima.

Crvene oči bile su uperene u njih.

Kralj Bili se ukočio.

U rukama je još držao zapaljene stranice.

``Idi!'' povikao je prema Šrajku. ``Vrati se odakle si došao!''

Šrajk se nije pomerio.

Silenus je uspeo da se pridigne.

``Gospodaru!'' viknuo je.

Kasnije nije znao kome se obraćao.

Biliju.

Ili Šrajku.

U jednom trenutku Kralj Bili je stajao pored fontane.

U sledećem je bio desetak metara dalje, podignut iznad zemlje.

Šrajk ga je držao.

Metalni prsti probili su mu ruke, grudi i noge.

Bili je još bio živ.

Još je držao deo zapaljenog Speva.

``Uništi ga!'' viknuo je.

Silenus nije znao na šta misli.

Na Šrajka?

Na rukopis?

Na oboje?

U fontani je ostalo više od hiljadu stranica.

Pored nje je bila kanta sa gorivom.

Silenus ju je uzeo.

Polio stranice.

Uzeo upaljač.

Pokušao da ga kresne.

Jednom.

Drugi put.

Treći.

Nije se upalio.

Pogledao je svoje životno delo.

Nije mogao.

Ispustio je upaljač.

Kralj Bili je vrištao dok su se Šrajkova sečiva pomerala kroz njegovo telo.

``Okončaj to, Martine!''

Silenus je podigao kantu.

I bacio je.

Ne na rukopis.

Na Bilija i Šrajka.

Gorivo ih je natopilo.

Plamen sa stranica koje je Bili još držao zahvatio je paru.

Obojica su nestala u vatri.

Silenus je video Bilija kako gori na Šrajkovim trnovima.

Čuo je njegov vrisak.

Onda su nestali.

Ne pali.

Ne izgoreli.

Jednostavno više nisu bili tamo.

Nije ostao pepeo.

Nije ostalo telo.

Samo poluspaljeni rukopis.

I Martin Silenus koji je klečao u dvorištu i vrištao.

\bigskip

Spev nije bio potpuno uništen.

Silenus je mesecima sastavljao preostale stranice.

Ponovo je pisao delove kojih se sećao.

Popravljao ono što je izgorelo.

Ali nije uspeo da nastavi tamo gde je stao.

Muza je nestala.

Grad Pesnika ostao je mrtav.

Silenus je još godinama lutao njegovim ruševinama.

Razgovarao sam sa sobom.

Sa Šrajkom.

Sa Vremenskim grobnicama.

Ponekad bi nailazili hodočasnici i videli dugokosog, bradatog luđaka kako iz ruševina psuje Boga bola i izaziva ga da izađe.

Na kraju je i to ludilo oslabilo.

Silenus je napustio grad.

Peške je prešao više od hiljadu kilometara do civilizacije.

U rancu je nosio samo rukopis.

\bigskip

Posle toga prošla su dva i po veka.

Tretmani za produženje života.

Duga putovanja.

Kriogeni san.

Godine koje je gubio dok je za ostatak čovečanstva prolazilo mnogo više vremena.

Čekao je.

I čuvao rukopis.

``Zašto?'' pitao je Sol Vejntraub.

Martin Silenus ga je pogledao kao da je odgovor očigledan.

``Zato što nije završen.''

``Posle toliko vremena još nameravaš da dovršiš poemu?''

``Naravno.''

Bron Lamija je prekrstila ruke.

``I zato ideš Šrajku? Zbog knjige?''

Silenusov uobičajeni osmeh nestao je.

Na trenutak je izgledao mnogo starije nego što je njegovo podmlađeno telo dopuštalo.

``Ne zbog knjige.''

Pogledao je prema severu.

Negde daleko ispred njih ležale su Vremenske grobnice.

``Zbog poslednjih stihova.''

Konzul se setio teškog prtljaga koji je pesnik nosio sa sobom.

``Rukopis je kod tebe?''

Silenus je klimnuo.

``Sve što je preživelo.''

``I očekuješ da će Šrajk ponovo postati tvoja muza?'' pitao je Hojt.

``Ne očekujem.''

Silenus je uzeo bocu.

Dugo je posmatrao vino pre nego što je otpio.

``Znam.''

Niko se nije odmah nasmejao.

Čak ni Silenus.

``Čitavog života pokušavao sam da napišem nešto što će ostati kada mene više ne bude'', rekao je. ``Bogatsvo nestane. Svetovi nestanu. Ljudi nestanu. Zemlja je nestala. Moja majka je nestala. Bili je nestao. Čitave civilizacije će nestati.''

Dodirnuo je prstom slepoočnicu.

``Ali reči ponekad ostanu.''

Pogledao je ostale.

``Možda je to jedina besmrtnost koju smo ikada stvarno imali.''

Ponovo je pogledao prema severu.

``Spev mora biti dovršen.''

Na trenutak niko nije govorio.

Onda se Martin Silenus ponovo nasmešio svojim bezobraznim, gotovo dečačkim osmehom.

``U početku je bila Reč.''

Podigao je bocu.

``A ako budem imao ikakve proklete sreće, na kraju će opet biti Reč.''